\documentclass[11pt]{article}
\usepackage[margin=1in]{geometry}
\usepackage{amsmath, amssymb}
\usepackage[shortlabels]{enumitem}
\usepackage{hyperref}
\usepackage{fancyhdr}
\usepackage{microtype}

\newcommand{\E}{\mathbb{E}}
\newcommand{\R}{\mathbb{R}}
\newcommand{\pd}[2]{\frac{\partial #1}{\partial #2}}

\newcounter{problem}
\newenvironment{problem}[1][]{%
  \stepcounter{problem}%
  \medskip
  \noindent\textbf{Problem \theproblem.}\ifx&#1&\else\quad(\textit{#1 pts})\fi\par\noindent\ignorespaces
}{\medskip}

\pagestyle{fancy}
\fancyhf{}
\lhead{ECON 320}
\chead{Problem Set 1}
\rhead{Due: [Date]}
\cfoot{\thepage}

\begin{document}

\begin{center}
  {\LARGE\bfseries ECON 320}\\[0.4em]
  {\large Problem Set 1}\\[0.2em]
  {\normalsize Due: [Date] \quad Instructor: [Name]}
\end{center}
\medskip\hrule\medskip

\noindent\textbf{Instructions:} Show all work. Partial credit is awarded for correct reasoning.

\vspace{1em}

\begin{problem}[10]
  Problem text here.
  \begin{enumerate}[(a)]
    \item
    \item
  \end{enumerate}
\end{problem}

\begin{problem}[15]
  Problem text here.
  \begin{enumerate}[(a)]
    \item
    \item
    \item
  \end{enumerate}
\end{problem}

\begin{problem}[10]
  Problem text here.
\end{problem}

\end{document}
